\documentclass[dvipdfmx, xcolor={dvipsnames},12pt]{beamer}
\usepackage{bxdpx-beamer}
\usepackage{pxjahyper}
\usepackage{otf}
\usetheme{CambridgeUS}

\usepackage[linguistics]{forest}
\usepackage{ascmac}
\usepackage{amsmath}
\usepackage{amsthm}
\usepackage{amssymb}
\usepackage{amsfonts}
\usepackage{latexsym}
\usepackage{mathtools}
\usepackage{float}
\usepackage{url}
\usepackage{xcolor}
\usepackage[most]{tcolorbox}
\usepackage{tikz}
\usepackage{varwidth}
\usetikzlibrary{positioning}

\newtcolorbox{marker}{
enhanced,
colback=yellow!20,
colframe=yellow!40!black,
boxrule=0.4pt,
left=8mm,
overlay={
% 左の!帯
\fill[brown!80!black] (frame.north west) rectangle ([xshift=6mm]frame.south west);
\node[white,font=\bfseries\Large] at ([xshift=3mm]frame.west) {!};
% 右下折れ曲がり
\fill[yellow!70!white]
(frame.south east) -- ++(-6mm,0) -- ++(6mm,6mm) -- cycle;
},
drop shadow,
}
\newtcolorbox{dfn}[2][]{%
enhanced,
title={#2},
fonttitle=\bfseries,
colback=cyan!5,
colframe=cyan!50!black,
sharp corners=northwest,
% --- ここからタイトルの見た目調整 ---
attach boxed title to top left={%
xshift=0mm, % 
yshift*=-0.5mm % 
},
varwidth boxed title*=-3mm,
boxed title style={%
colback=cyan!50!black,
colframe=cyan!50!black,
rounded corners,
sharp corners=south,
boxrule=0.5pt,
},
#1
}


\newtcolorbox{thm}[2][]{%
enhanced,
title={#2},
fonttitle=\bfseries,
colback=green!5,
colframe=green!50!black,
sharp corners=northwest,
% --- ここからタイトルの見た目調整 ---
attach boxed title to top left={%
xshift=0mm, % 左端から少し内側
yshift*=-0.5mm % 枠線に半分かぶせる
},
varwidth boxed title*=-3mm, % タイトルの幅に合わせた小さい箱
boxed title style={%
colback=green!50!black,
colframe=green!50!black,
rounded corners,
sharp corners=south,
boxrule=0.5pt,
},
#1
}



\title{Chomsky (1956) Introduction }
\author{荒木 理求}
\date\today


\begin{document}

\frame{\titlepage}

\begin{frame}{二つの中心的問題}
Chomsky (1956)

言語の記述的な研究 (the descriptive study of language) には2つの中心的問題がある :

\begin{enumerate}
    \item 各自然言語に対して\\
    \quad \textbf{単純で「本質を明らかにする」文法を発見すること}
    
    \item それらの文法の性質を分析し,\\
    \quad \textbf{一般的な言語理論へ到達すること}
\end{enumerate}

\begin{marker}
記述言語学 (descriptive linguistics) との違いに注意せよ. 
\end{marker}

\begin{marker}
一般理論の構築には個別言語の分析が必須\\
(cf. Ferdinand de Saussure)
\end{marker}
\end{frame}

%--------------------------------------------------
\begin{frame}{言語理論の基本的視点}
\begin{itemize}
    \item 文法は言語の構造に関する理論と見ることができる.
    \item 理論の役割：
    \begin{enumerate}
        \item 有限の観察（コーパス）から
        \item 無限個の文の集合を特徴づける
    \end{enumerate}
\end{itemize}
    標語的に言えば
    \begin{block}{}
        文法は観察された文の集合を「正文の無限集合」へ射影する
    \end{block}
\begin{marker}
数学では, Euclid原論の時代から, 公理 (命題の集合) から出発して定理 (嬉しい結果) を得ることを目標にしている. 1920年代以降は, ヒルベルト$\cdot$プログラムやゲーデルの活躍によって, 数学のさらなる形式化が進んだ.
\end{marker}
\end{frame}

%--------------------------------------------------
\begin{frame}{文法性の直観（例）}
英語話者は次を区別できる：

\begin{enumerate}
\item[(1)] John ate a sandwich
\item[(2)] Sandwich ate John
\end{enumerate}

\begin{itemize}
    \item (1)：文法的
    \item (2)：非文法的
\end{itemize}
文法はこの区別を正しく再現する必要がある

目標とすべきモデルとしての文法：
\begin{itemize}
    \item (1) を含み
    \item (2) を排除する
\end{itemize}
\end{frame}

%--------------------------------------------------
\begin{frame}{言語理論の目的}
言語理論（general linguistic theory）とは：

\begin{block}{}
有限のコーパスから、\\
どの文法を選ぶべきかを決定するためのメタ理論
\end{block}

話者の能力, すなわち
    \begin{itemize}
        \item 新しい文を生成できる
        \item 非文法文を排除できる
    \end{itemize}
能力を, 制限された経験から身に着けるメカニズムを説明したい.
\end{frame}

%--------------------------------------------------
\begin{frame}{文法の本質}
\begin{itemize}
    \item 文法は有限な観察に基づく
    \item しかし生成する対象は無限個の文
\end{itemize}
そのために一般的規則（rules）を導入


適切な文法は
文法的文の集合を一意に定めるべき
\end{frame}

%--------------------------------------------------
\begin{frame}{言語の形式的定義}

\begin{dfn}{定義1$^{\text{lx}}$(語, 言語)}
$A$をアルファベットからなる有限集合として
\begin{itemize}
\item $A$上の語 : Aのアルファベットをつなげた (concatenating) 記号列（string）
\item 言語 $L$：有限の長さの語の集合
\item 文法 : $L$ の文をすべて、かつそれのみ生成する装置
\end{itemize}
\end{dfn}

\end{frame}

%--------------------------------------------------
\begin{frame}{理論評価の基本方針}
\begin{itemize}
    \item 明確な事例でテストする
    \item 文法は
    \begin{enumerate}
        \item 正文を含み
        \item 非文を排除する
    \end{enumerate}
　　する必要がある
\end{itemize}

これは個別言語だけではなく, 必要条件として, 理論全体で力を発揮するテストになる.

\end{frame}

%--------------------------------------------------
\begin{frame}{重要な問い}
\textbf{(3)}\\
\begin{block}{}
提案された型の理論では記述できないような興味深い言語は存在するか？
\end{block}

\begin{itemize}
    \item YES → 理論は不十分
    \item NO → 次の問いへ
\end{itemize}
\end{frame}

%--------------------------------------------------
\begin{frame}{重要な問い (2)}
\textbf{(4)}\\
\begin{block}{}
すべての興味深い言語に対して十分に単純な文法を構成できるか？
\end{block}
\end{frame}

%--------------------------------------------------
\begin{frame}{重要な問い}
\textbf{(5)}\\
\begin{block}{}
その文法は「明らかにする (revealing)」ものか？\\
\end{block}
すなわち
\begin{itemize}
    \item 統語構造が意味の分析を支えるか
    \item 言語の使用・理解に洞察を与えるか
\end{itemize}
\end{frame}

\begin{frame}
\frametitle{Appendix}
\begin{dfn}{定義2. 句構造文法.}
次のシステム$G\coloneq\langle V_\text{N}, V_\text{T}, S, P\rangle$を\textbf{句構造文法}という.
\begin{itemize}
\item $V_\text{N}$はアルファベット, その元を非終端記号 (nonterminal symbol) または変数という.
\item $V_\text{T}$もアルファベットで $V_\text{N}\cap V_\text{T}=\emptyset$を満たす. またその元を終端記号 (terminal symbol) という.
\item $S\in V_\text{N}$で, 開始記号という.
\item $P$は$u\to v$という形の式の集合である. ただし$u,v\in (V_\text{N}\cup V_\text{T})^*$ (すなわち$V_\text{N}\cup V_\text{T}$上の語全体の集合) で, かつ$u$は少なくとも一つ変数を含む. $P$の元を生成規則という.
\end{itemize}
\end{dfn}
\end{frame}

\begin{frame}
\frametitle{Appendix}
句構造文法は単に文法ともよぶ. さてこれで文脈自由文法が定義できる.

\begin{dfn}{定義3. 文脈自由文法.}
文法$G=\langle V_\text{N}, V_\text{T}, S, P\rangle$, $V\coloneq V_\text{N}\cup V_\text{T}$において, $P$の生成規則が
\[
A\to \beta \hspace{20pt}(A\in V_\text{N}, \beta\in V^*)
\]
という形の物だけであるとき\footnote{正確には, 開始記号$S$以外の右辺には空語$\lambda$が現れないという制約をつける.}, $G$は\textbf{文脈自由文法 (Context-Free Grammar; CFG)}であるといい, $G$によって生成される言語を\textbf{文脈自由言語(Context-Free Language; CFL)}という.
\end{dfn}
\end{frame}



\end{document}